\documentclass[a4paper]{article}

\usepackage[english]{babel}
%\usepackage[utf8x]{inputenc}
\usepackage{amsmath}
\usepackage{graphicx}
\usepackage[colorinlistoftodos]{todonotes}
\usepackage[all]{xy}
\usepackage{dsfont}
\usepackage{amssymb}
\usepackage{bbold}
\usepackage{amsthm}

\usepackage{tikz}
\usepackage{verbatim}
\usetikzlibrary{matrix}
\usetikzlibrary{cd}

\usepackage[autostyle=true]{csquotes}
\usepackage{pdfpages}
\usepackage{hyperref}
\usepackage{fancyhdr}



\newtheorem{theorem}{Theorem}[section]
\newtheorem{proposition}{Proposition}[theorem]
\newtheorem{lemma}{Lemma}[theorem]
\newtheorem{corollary}{Corollary}[theorem]
\newtheorem*{definition}{Definition}


\newcommand{\norm}[1]{\left\lVert#1\right\rVert}

\newcommand{\Rn}{\mathbb{R}^n}


\newcommand{\too}[1]{\overset{#1}{\longrightarrow}}
\newcommand{\hookar}[1]{\overset{#1}{\hookrightarrow}}
\newcommand{\colim}[1]{ \underset{#1}{\texttt{colim}} \text{ } }
\newcommand{\fset}[1]{ \{ #1_0 , \dots , #1_n \} }
\newcommand{\ddt}{\frac{d}{dt}}


\newcommand{\N}{\mathbb{N}}
\newcommand{\Z}{\mathbb{Z}}







\fancyhead[L]{\includegraphics[scale=0.1]{UFT.png}}

\begin{document}


\thispagestyle{fancy}
\vspace*{1cm}




\section*{Preguntas}
\begin{enumerate}


    \item Considere la siguiente matriz
\[    A=\begin{pmatrix} 1 &0 & -1 \\ 0 & -1 & 2 \\ 0 & 0 & 2 \end{pmatrix}.\]
\begin{enumerate}
    \item Encuentre los valores propios de $A$.
    \item Determine si es diagonalizable, en caso positivo, encuentre la matriz diagonal.
\end{enumerate}
    \item Considere la siguiente matriz
\[    A=\begin{pmatrix} 1 &2 & 1 \\ 0 & -1 & -2 \\ 0 & 0 & 3 \end{pmatrix}.\]
\begin{enumerate}
    \item Encuentre los valores propios de $A$.
    \item Determine si es diagonalizable, en caso positivo, encuentre la matriz diagonal.
\end{enumerate}


        









\end{enumerate}


\end{document}